% c18 原书 p.223–235（PDF p.230–242）
\subsection{Grothendieck 存在定理}\label{sec:2.7}

设 $A$ 是诺特环，$I$ 是 $A$ 的理想，$X$ 是 $\Spec A$ 上的固有概形，$\cF$ 是凝聚 $\cO_X$-模，$\widehat{X}$ 与 $\widehat{\cF}$ 分别是 $X$ 与 $\cF$ 关于 $I$ 的形式完全化。本节我们证明
\begin{equation*}
H^n(\widehat{X},\widehat{\cF})\cong\plim_k H^n(X,\cF/I^k\cF).
\end{equation*}
此外我们还证明 Grothendieck 存在定理，它断言每个凝聚 $\cO_{\widehat{X}}$-模都是某个凝聚 $\cO_X$-模的形式完全化。先从一系列引理开始。

\begin{lemma}\label{it:2.7.1}
设 $(C_k^{\bullet})_{k\in\bN}$ 是阿贝尔群复形的逆向系统。假设对每个 $n$，逆向系统 $(H^n(C_k^{\bullet}))_{k\in\bN}$ 满足 Mittag-Leffler 条件，且对每个 $n$ 与每对 $l\ge k$，同态 $C_l^n\to C_k^n$ 是满射。则对每个 $n$，我们有
\begin{equation*}
H^n(\plim_k C_k^{\bullet})\cong\plim_k H^n(C_k^{\bullet}).
\end{equation*}
\end{lemma}

\begin{proof}
对每对 $l\ge k$ 与每个 $n$，由于 $C_l^{n-1}\to C_k^{n-1}$ 是满射，$B^n(C_l^{\bullet})\to B^n(C_k^{\bullet})$ 也是满射。由于 $(H^n(C_k^{\bullet}))_{k\in\bN}$ 满足 Mittag-Leffler 条件，存在依赖于 $k$ 的 $k_0$，使得对每个 $l\ge k+k_0$，有
\begin{equation*}
\Image\big(H^n(C_l^{\bullet})\to H^n(C_k^{\bullet})\big)=\Image\big(H^n(C_{k+k_0}^{\bullet})\to H^n(C_k^{\bullet})\big).
\end{equation*}
我们有交换图

\begin{center}
\begin{tikzcd}
0 \arrow[r] & B^n(C_l^{\bullet}) \arrow[r] \arrow[d] & Z^n(C_l^{\bullet}) \arrow[r] \arrow[d] & H^n(C_l^{\bullet}) \arrow[r] \arrow[d] & 0 \\
0 \arrow[r] & B^n(C_k^{\bullet}) \arrow[r] & Z^n(C_k^{\bullet}) \arrow[r] & H^n(C_k^{\bullet}) \arrow[r] & 0.
\end{tikzcd}
\end{center}

由追图可证，对每个 $l\ge k+k_0$，有
\begin{equation*}
\Image\big(Z^n(C_l^{\bullet})\to Z^n(C_k^{\bullet})\big)=\Image\big(Z^n(C_{k+k_0}^{\bullet})\to Z^n(C_k^{\bullet})\big).
\end{equation*}
所以逆向系统 $(Z^n(C_k^{\bullet}))_{k\in\bN}$、$(B^n(C_k^{\bullet}))_{k\in\bN}$ 与 $(H^n(C_k^{\bullet}))_{k\in\bN}$ 对一切 $n$ 都满足 Mittag-Leffler 条件。由 1.5.1，下列序列正合：
\begin{align*}
0\to\plim_k Z^n(C_k^{\bullet})&\to\plim_k C_k^n\to\plim_k B^{n+1}(C_k^{\bullet})\to 0,\\
0\to\plim_k B^n(C_k^{\bullet})&\to\plim_k Z^n(C_k^{\bullet})\to\plim_k H^n(C_k^{\bullet})\to 0.
\end{align*}
这说明 $H^n(\plim_k C_k^{\bullet})\cong\plim_k H^n(C_k^{\bullet})$。
\end{proof}

\begin{tnote}
引理 2.7.1 是本节反复使用的技术核心：Mittag-Leffler 条件保证上同调与逆向极限可以交换，而「每对 $l\ge k$ 同态满射」则保证圈层与边界层的逆向系统也具有所需性质。原书随后在 2.7.2 与 2.7.3 中把这一交换性从复形逐层搬到层与上同调，构成 2.7.4 证明的骨架。[依据：原书 p.223–226]
\end{tnote}

\begin{lemma}\label{it:2.7.2}
设 $(\cF_k)_{k\in\bN}$ 是拓扑空间 $X$ 上阿贝尔群层的逆向系统。假设每个 $\cF_k$ 是松弛层，且对每对 $l\ge k$，态射 $\cF_l\to\cF_k$ 是满射且其核是松弛层。则对 $X$ 的每个开子集 $U$，同态
\begin{equation*}
\Gamma(U,\cF_l)\to\Gamma(U,\cF_k)
\end{equation*}
对每对 $l\ge k$ 都是满射，且对每个 $n\ge1$ 有
\begin{equation*}
H^n(U,\plim_k\cF_k)=0.
\end{equation*}
\end{lemma}

\begin{proof}
设 $\cK_{lk}$ 是满射 $\cF_l\to\cF_k$ 的核。序列
\begin{equation*}
0\to\cK_{lk}\to\cF_l\to\cF_k\to 0
\end{equation*}
是松弛层的正合序列。由引理 2.1.11，对 $X$ 的每个开子集 $U$，
\begin{equation*}
\Gamma(U,\cF_l)\to\Gamma(U,\cF_k)
\end{equation*}
是满射。此外，对 $U$ 的每个开覆盖 $\mathfrak{U}$ 与每个 $n$，同态
\begin{equation*}
C^n(\mathfrak{U},\cF_l)\to C^n(\mathfrak{U},\cF_k)
\end{equation*}
是满射。注意对每个 $n$，
\begin{equation*}
H^n(C^{\bullet}(\mathfrak{U},\cF_l))\to H^n(C^{\bullet}(\mathfrak{U},\cF_k))
\end{equation*}
是满射。事实上，当 $n=0$ 时，这由 $\Gamma(U,\cF_l)\to\Gamma(U,\cF_k)$ 的满射性得到。对每个 $n\ge1$，由 2.3.6 我们有
\begin{equation*}
H^n(C^{\bullet}(\mathfrak{U},\cF_k))=0.
\end{equation*}
于是我们可以对逆向系统 $(C^{\bullet}(\mathfrak{U},\cF_k))_{k\in\bN}$ 应用 2.7.1，得到
\begin{equation*}
H^n\big(\plim_k C^{\bullet}(\mathfrak{U},\cF_k)\big)\cong\plim_k H^n(C^{\bullet}(\mathfrak{U},\cF_k)).
\end{equation*}
注意
\begin{equation*}
\plim_k C^{\bullet}(\mathfrak{U},\cF_k)=C^{\bullet}\big(\mathfrak{U},\plim_k\cF_k\big).
\end{equation*}
所以我们有
\begin{equation*}
\Cech^n\big(\mathfrak{U},\plim_k\cF_k\big)\cong\plim_k\Cech^n(\mathfrak{U},\cF_k).
\end{equation*}
当 $n\ge1$ 时，由 2.3.6 有 $H^n(\mathfrak{U},\cF_k)=0$，从而
\begin{equation*}
\Cech^n\big(\mathfrak{U},\plim_k\cF_k\big)=0.
\end{equation*}
这对每个开集 $U$ 的每个开覆盖 $\mathfrak{U}$ 都成立。由 2.3.11，这蕴含对每个 $n\ge1$ 有
\begin{equation*}
H^n\big(X,\plim_k\cF_k\big)=0.
\end{equation*}
\end{proof}

\begin{lemma}\label{it:2.7.3}
设 $X$ 是拓扑空间，$\mathfrak{B}$ 是 $X$ 的一族开子集，它构成 $X$ 的拓扑的基，并设 $(\cF_k)_{k\in\bN}$ 是 $X$ 上阿贝尔群层的逆向系统。假设下列条件成立：
\begin{enumerate}[label=(\alph*),leftmargin=2.8em]
\item 对 $\mathfrak{B}$ 中的每个开子集 $U$ 与每对 $l\ge k$，同态
\begin{equation*}
\Gamma(U,\cF_l)\to\Gamma(U,\cF_k)
\end{equation*}
是满射。
\item 对 $\mathfrak{B}$ 中的每个 $U$，每个 $k$ 与每个 $n\ge1$，有
\begin{equation*}
H^n(U,\cF_k)=0.
\end{equation*}
\item 对每个 $n$，逆向系统 $(H^n(X,\cF_k))_{k\in\bN}$ 满足 Mittag-Leffler 条件。
\end{enumerate}
则我们有
\begin{equation*}
H^n\big(X,\plim_k\cF_k\big)\cong\plim_k H^n(X,\cF_k).
\end{equation*}
\end{lemma}

\begin{proof}
由 (a)，对每对 $l\ge k$，态射 $\cF_l\to\cF_k$ 是满射。设 $\cK_{lk}$ 是它的核。由 2.1.10，对每个 $i\ge0$ 我们有正合序列
\begin{equation*}
0\to C^i(\cK_{lk})\to C^i(\cF_l)\to C^i(\cF_k)\to 0,
\end{equation*}
其中 $C^{\bullet}(-)$ 是 Godement 分解函子。于是对每个 $i$，逆向系统 $(C^i(\cF_k))_{k\in\bN}$ 满足 2.7.2 的条件。因此对 $X$ 的每个开子集 $U$ 与每对 $l\ge k$，同态
\begin{equation*}
\Gamma(U,C^i(\cF_l))\to\Gamma(U,C^i(\cF_k))
\end{equation*}
是满射，且对每个 $n\ge1$ 与每个 $i$ 有
\begin{equation*}
H^n\big(U,\plim_k C^i(\cF_k)\big)=0.
\end{equation*}
设 $U$ 是 $\mathfrak{B}$ 中的一个成员。由 (a)、(b) 以及
\begin{equation*}
H^n(\Gamma(U,C^{\bullet}(\cF_k)))\cong H^n(U,\cF_k),
\end{equation*}
可知对每个 $n$ 与每对 $l\ge k$，同态
\begin{equation*}
H^n(\Gamma(U,C^{\bullet}(\cF_l)))\to H^n(\Gamma(U,C^{\bullet}(\cF_k)))
\end{equation*}
是满射。我们已经看到
\begin{equation*}
\Gamma(U,C^n(\cF_l))\to\Gamma(U,C^n(\cF_k))
\end{equation*}
是满射。于是由 2.7.1 得到
\begin{equation*}
H^n\big(\plim_k\Gamma(U,C^{\bullet}(\cF_k))\big)\cong\plim_k H^n(\Gamma(U,C^{\bullet}(\cF_k)))=\plim_k H^n(U,\cF_k).
\end{equation*}
当 $n\ge1$ 时，由 (b) 有
\begin{equation*}
H^n\big(\plim_k\Gamma(U,C^{\bullet}(\cF_k))\big)\cong\plim_k H^n(U,\cF_k)=0.
\end{equation*}
这对族 $\mathfrak{B}$ 中的每个 $U$ 都成立。所以序列
\begin{equation*}
0\to\plim_k\cF_k\to\plim_k C^0(\cF_k)\to\plim_k C^1(\cF_k)\to\cdots
\end{equation*}
正合。我们已经看到，对每个 $n\ge1$ 与每个 $i$ 有
\begin{equation*}
H^n\big(X,\plim_k C^i(\cF_k)\big)=0.
\end{equation*}
由 2.1.2，我们有
\begin{equation*}
H^n\big(X,\plim_k\cF_k\big)\cong H^n\big(\Gamma(X,\plim_k C^{\bullet}(\cF_k))\big)=H^n\big(\plim_k\Gamma(X,C^{\bullet}(\cF_k))\big).
\end{equation*}
我们已经看到 $\Gamma(X,C^n(\cF_l))\to\Gamma(X,C^n(\cF_k))$ 是满射。此外，对每个 $n$，逆向系统
\begin{equation*}
(H^n(\Gamma(X,C^{\bullet}(\cF_k))))_{k\in\bN}=(H^n(X,\cF_k))_{k\in\bN}
\end{equation*}
由 (c) 满足 Mittag-Leffler 条件。于是再次由 2.7.1 得到
\begin{equation*}
H^n\big(\plim_k\Gamma(X,C^{\bullet}(\cF_k))\big)\cong\plim_k H^n(\Gamma(X,C^{\bullet}(\cF_k)))\cong\plim_k H^n(X,\cF_k).
\end{equation*}
因此
\begin{equation*}
H^n\big(X,\plim_k\cF_k\big)\cong\plim_k H^n(X,\cF_k).
\end{equation*}
\end{proof}

设 $X$ 是诺特概形，$\cI$ 是 $\cO_X$ 的凝聚理想层。回忆 $X$ 关于 $\cI$ 的形式完全化是由
\begin{align*}
\widehat{X}&=\Supp(\cO_X/\cI),\\
\cO_{\widehat{X}}&=\plim_k(\cO_X/\cI^k)|_{\widehat{X}}
\end{align*}
定义的局部环化空间 $(\widehat{X},\cO_{\widehat{X}})$。对每个 $\cO_X$-模 $\cF$，它的形式完全化是 $\cO_{\widehat{X}}$-模
\begin{equation*}
\widehat{\cF}=\plim_k(\cF/\cI^k\cF)|_{\widehat{X}}.
\end{equation*}

\begin{theorem}\label{it:2.7.4}
设 $A$ 是诺特环，$I$ 是 $A$ 的理想，$X$ 是 $\Spec A$ 上的固有概形，$\cF$ 是凝聚 $\cO_X$-模，$\widehat{X}$ 与 $\widehat{\cF}$ 分别是 $X$ 与 $\cF$ 关于理想层 $I^{\sim}\cO_X$ 的形式完全化。则对每个 $n$，我们有同构
\begin{equation*}
H^n(\widehat{X},\widehat{\cF})\cong\plim_k H^n(X,\cF/I^k\cF),
\end{equation*}
其中 $I^k\cF=(I^k)^{\sim}\cF$。
\end{theorem}

\begin{proof}
设 $i\colon\widehat{X}\to X$ 是局部环化空间的标准态射。由于每个 $\cF/I^k\cF$ 都支撑在闭子集 $\widehat{X}$ 上，我们有
\begin{equation*}
\cF/I^k\cF=i_*\big((\cF/I^k\cF)|_{\widehat{X}}\big),
\end{equation*}
从而
\begin{equation*}
i_*\widehat{\cF}=i_*\big(\plim_k(\cF/I^k\cF)|_{\widehat{X}}\big)=\plim_k\cF/I^k\cF.
\end{equation*}
利用 $i_*$ 是正合函子、且把 $\widehat{X}$ 上的内射层变为 $X$ 上的内射层这一事实，可以验证
\begin{equation*}
H^n(\widehat{X},\widehat{\cF})\cong H^n(X,i_*\widehat{\cF}).
\end{equation*}
所以
\begin{equation*}
H^n(\widehat{X},\widehat{\cF})\cong H^n\big(X,\plim_k\cF/I^k\cF\big).
\end{equation*}
我们有长正合序列
\begin{equation*}
\cdots\to H^n(X,\cF)\to H^n(X,\cF/I^k\cF)\to H^{n+1}(X,I^k\cF)\to\cdots.
\end{equation*}
设 $Q_k=\Image\big(H^n(X,\cF/I^k\cF)\to H^{n+1}(X,I^k\cF)\big)$。则我们有正合序列
\begin{equation*}
H^n(X,\cF)\to H^n(X,\cF/I^k\cF)\to Q_k\to 0.
\end{equation*}
由 2.5.8，存在自然数 $k_0$，使得对每个 $k\ge k_0$，标准同态 $Q_{k+k_0}\to Q_k$ 为零。于是对每个 $k$，同态 $Q_{k+2k_0}\to Q_{k+k_0}$ 也为零。因为对每个 $k$ 与每个 $l\ge k+2k_0$，标准同态 $Q_l\to Q_k$ 可以分解为
\begin{equation*}
Q_l\to Q_{k+2k_0}\to Q_{k+k_0}\to Q_k,
\end{equation*}
所以它为零。
考虑如下交换图（$l\ge k+2k_0$）：

\begin{center}
\begin{tikzcd}[column sep=small]
H^n(X,\cF) \arrow[r] \arrow[d,equal] & H^n(X,\cF/I^l\cF) \arrow[r] \arrow[d] & Q_l \arrow[r] \arrow[d] & 0 \\
H^n(X,\cF) \arrow[r] & H^n(X,\cF/I^k\cF) \arrow[r] & Q_k \arrow[r] & 0
\end{tikzcd}
\end{center}

由于第三个竖直箭头为零，我们有
\begin{equation*}
\Image\big(H^n(X,\cF/I^l\cF)\to H^n(X,\cF/I^k\cF)\big)\subset\Image\big(H^n(X,\cF)\to H^n(X,\cF/I^k\cF)\big).
\end{equation*}
由图中左方交换方块，我们有
\begin{equation*}
\Image\big(H^n(X,\cF)\to H^n(X,\cF/I^k\cF)\big)\subset\Image\big(H^n(X,\cF/I^l\cF)\to H^n(X,\cF/I^k\cF)\big).
\end{equation*}
所以对每个 $l\ge k+2k_0$，
\begin{equation*}
\Image\big(H^n(X,\cF/I^l\cF)\to H^n(X,\cF/I^k\cF)\big)=\Image\big(H^n(X,\cF)\to H^n(X,\cF/I^k\cF)\big).
\end{equation*}
于是逆向系统 $(H^n(X,\cF/I^k\cF))_{k\in\bN}$ 满足 Mittag-Leffler 条件。另一方面，对 $X$ 的每个仿射开子概形 $U$，同态
\begin{equation*}
\Gamma(U,\cF/I^l\cF)\to\Gamma(U,\cF/I^k\cF)
\end{equation*}
对每对 $l\ge k$ 都是满射，且对每个 $k$ 与每个 $n\ge1$ 有
\begin{equation*}
H^n(U,\cF/I^k\cF)=0.
\end{equation*}
由 2.7.3，我们有
\begin{equation*}
H^n\big(X,\plim_k\cF/I^k\cF\big)\cong\plim_k H^n(X,\cF/I^k\cF).
\end{equation*}
所以
\begin{equation*}
H^n(\widehat{X},\widehat{\cF})\cong\plim_k H^n(X,\cF/I^k\cF).
\end{equation*}
\end{proof}

\begin{corollary}\label{it:2.7.5}
设 $A$ 是诺特环，$I$ 是 $A$ 的理想，$X$ 是 $\Spec A$ 上的固有概形，$\cF$ 是凝聚 $\cO_X$-模，$\widehat{X}$ 与 $\widehat{\cF}$ 分别是 $X$ 与 $\cF$ 关于理想层 $I^{\sim}\cO_X$ 的形式完全化。则对每个 $n$，我们有同构
\begin{equation*}
H^n(\widehat{X},\widehat{\cF})\cong(H^n(X,\cF))^{\wedge},
\end{equation*}
其中 $(H^n(X,\cF))^{\wedge}$ 是 $H^n(X,\cF)$ 的 $I$-进完全化。
\end{corollary}

\begin{proof}
由 2.5.6 与 2.7.4 即得。
\end{proof}

\begin{corollary}\label{it:2.7.6}
设 $f\colon X\to Y$ 是诺特概形的固有态射，$\cI$ 是 $\cO_Y$ 的凝聚理想层，$\cF$ 是凝聚 $\cO_X$-模，$X_k$ 与 $Y_k$ 分别是 $X$ 与 $Y$ 的由理想层 $\cI^k\cO_X$ 与 $\cI^k$ 定义的闭子概形，$\cF_k=(\cF/I^k\cF)|_{X_k}$，$\widehat{X}$、$\widehat{\cF}$ 与 $\widehat{Y}$ 分别是 $X$、$\cF$ 与 $Y$ 关于 $\cI\cO_X$、$\cI\cO_X$ 与 $\cI$ 的形式完全化，$f_k\colon X_k\to Y_k$ 与 $\hat{f}\colon\widehat{X}\to\widehat{Y}$ 是由 $f$ 诱导的态射。则对每个 $n$，我们有同构
\begin{equation*}
R^n\hat{f}_*\widehat{\cF}\cong\plim_k R^nf_{k*}\cF_k\cong\widehat{R^nf_*\cF},
\end{equation*}
其中 $\widehat{R^nf_*\cF}$ 是 $R^nf_*\cF$ 关于 $\cI$ 的形式完全化。
\end{corollary}

\begin{proof}
注意 $\widehat{X}$ 与 $X_k$（resp.\ $\widehat{Y}$ 与 $Y_k$）有相同的底拓扑空间，且 $\hat{f}$ 与 $f_k$ 作为这些拓扑空间上的映射相同。标准态射族 $\widehat{\cF}\to\cF_k$ 诱导态射族
\begin{equation*}
R^n\hat{f}_*\widehat{\cF}\to R^nf_{k*}\cF_k,
\end{equation*}
从而诱导态射
\begin{equation*}
R^n\hat{f}_*\widehat{\cF}\to\plim_k R^nf_{k*}\cF_k.
\end{equation*}
我们来证明它是同构。设 $U=\Spec A$ 是 $Y$ 的仿射开子概形。由 2.7.4 我们有
\begin{equation*}
H^n(\widehat{X}\cap f^{-1}(U),\widehat{\cF})\cong\plim_k H^n(f^{-1}(U),\cF/I^k\cF).
\end{equation*}
设 $i_{X_k}\colon X_k\to X$ 是闭浸入。我们有 $i_{X_k*}\cF_k=\cF/I^k\cF$。由于 $i_{X_k*}$ 是正合函子且把内射层变为内射层，我们有
\begin{equation*}
H^n(f^{-1}(U),\cF/I^k\cF)\cong H^n(X_k\cap f^{-1}(U),\cF_k).
\end{equation*}
所以
\begin{equation*}
H^n(\widehat{X}\cap f^{-1}(U),\widehat{\cF})\cong\plim_k H^n(X_k\cap f^{-1}(U),\cF_k),
\end{equation*}
即
\begin{equation*}
H^n(\hat{f}^{-1}(\widehat{Y}\cap U),\widehat{\cF})\cong\plim_k H^n(f_k^{-1}(Y_k\cap U),\cF_k).
\end{equation*}
但 $R^n\hat{f}_*\widehat{\cF}$ 是与前层
\begin{equation*}
V\mapsto H^n(\hat{f}^{-1}(V),\widehat{\cF})
\end{equation*}
相伴的层，且由 2.4.9 我们有
\begin{equation*}
(R^nf_{k*}\cF_k)(Y_k\cap U)\cong H^n(f_k^{-1}(Y_k\cap U),\cF_k).
\end{equation*}
当 $U$ 遍历 $Y$ 的仿射开子概形族时，$\widehat{Y}\cap U=Y_k\cap U$ 构成 $\widehat{Y}=Y_k$ 的拓扑的基。故我们有
\begin{equation*}
R^n\hat{f}_*\widehat{\cF}\cong\plim_k R^nf_{k*}\cF_k.
\end{equation*}
设 $i_{Y_k}\colon Y_k\to Y$ 是闭浸入。我们有 $fi_{X_k}=i_{Y_k}f_k$，从而
\begin{equation*}
R^n(fi_{X_k})_*\cF_k\cong R^n(i_{Y_k}f_k)_*\cF_k.
\end{equation*}
利用 $i_{X_k*}$ 与 $i_{Y_k*}$ 是正合函子且把内射层变为内射层这一事实，可以验证
\begin{align*}
R^n(fi_{X_k})_*\cF_k&\cong R^nf_*(i_{X_k*}\cF_k)\cong R^nf_*(\cF/I^k\cF),\\
R^n(i_{Y_k}f_k)_*\cF_k&\cong i_{Y_k*}R^nf_{k*}\cF_k.
\end{align*}
所以
\begin{equation*}
R^nf_*(\cF/I^k\cF)\cong i_{Y_k*}R^nf_{k*}\cF_k.
\end{equation*}
由 2.5.9 我们有
\begin{equation*}
\plim_k R^nf_*\cF/I^kR^nf_*\cF\cong\plim_k R^nf_*(\cF/I^k\cF).
\end{equation*}
所以
\begin{equation*}
\plim_k R^nf_*\cF/I^kR^nf_*\cF\cong\plim_k i_{Y_k*}R^nf_{k*}\cF_k.
\end{equation*}
结合两边逆向系统中的层都支撑在 $\widehat{Y}$ 上这一事实，这蕴含
\begin{equation*}
\plim_k(R^nf_*\cF/I^kR^nf_*\cF)|_{\widehat{Y}}\cong\plim_k R^nf_{k*}\cF_k,
\end{equation*}
即
\begin{equation*}
\widehat{R^nf_*\cF}\cong\plim_k R^nf_{k*}\cF_k.
\end{equation*}
\end{proof}

\begin{lemma}\label{it:2.7.7}
设 $X$ 是诺特概形，$\cI$ 是 $\cO_X$ 的凝聚理想层，$\widehat{X}$ 是 $X$ 关于 $\cI$ 的形式完全化，$i\colon\widehat{X}\to X$ 是标准态射。则函子 $\cF\mapsto i^*\cF$ 在 $\cO_X$-模范畴上正合。
\end{lemma}

\begin{proof}
由 1.5.13 (i) 与 (ii)，函子 $\cF\mapsto i^*\cF$ 在凝聚 $\cO_X$-模范畴上正合。一般地，对每个 $\cO_X$-模 $\cF$ 与每个 $P\in\widehat{X}$，我们有
\begin{equation*}
(i^*\cF)_P\cong\cO_{\widehat{X},P}\otimes_{\cO_{X,P}}\cF_P.
\end{equation*}
要证明 $i^*$ 在 $\cO_X$-模范畴上正合，只需证明对每个 $P\in\widehat{X}$，$\cO_{\widehat{X},P}$ 是平坦 $\cO_{X,P}$-代数。对每个有限生成 $\cO_{X,P}$-模 $M$，我们可以找到正合序列
\begin{equation*}
\cO_{X,P}^m\to\cO_{X,P}^n\to M\to 0
\end{equation*}
（对某些自然数 $m$ 与 $n$）。由 1.4.1 (i)，存在 $P$ 在 $X$ 中的开邻域 $U$ 与态射 $\cO_U^m\to\cO_U^n$，它在茎上诱导同态 $\cO_{X,P}^m\to\cO_{X,P}^n$。设 $\cF$ 是 $\cO_U^m\to\cO_U^n$ 的余核。则 $\cF$ 是凝聚 $\cO_U$-模，且 $\cF_P\cong M$。设
\begin{equation*}
0\to M'\to M\to M''\to 0
\end{equation*}
是有限生成 $\cO_{X,P}$-模的正合序列。我们可以找到 $P$ 在 $X$ 中的邻域 $U$ 与凝聚 $\cO_U$-模 $\cF'$、$\cF$、$\cF''$，使得它们在 $P$ 处的茎分别同构于 $M'$、$M$、$M''$。必要时缩小 $U$，我们可以假设存在态射 $\cF'\to\cF$ 与 $\cF\to\cF''$，它们在茎上分别诱导 $M'\to M$ 与 $M\to M''$。由于 $M'\to M$ 是单射，$\ker(\cF'\to\cF)$ 在 $P$ 处的茎为 $0$。由 1.4.1 (ii)，缩小 $U$ 后我们有 $\ker(\cF'\to\cF)=0$。类似地，缩小 $U$ 后我们得到凝聚 $\cO_U$-模的正合序列
\begin{equation*}
0\to\cF'\to\cF\to\cF''\to 0,
\end{equation*}
它在茎上诱导正合序列
\begin{equation*}
0\to M'\to M\to M''\to 0.
\end{equation*}
由 1.5.13 (i) 与 (ii)，序列
\begin{equation*}
0\to(i^*\cF')_P\to(i^*\cF)_P\to(i^*\cF'')_P\to 0
\end{equation*}
正合，其中 $i$ 是从 $U$ 的形式完全化 $\widehat{U}$ 到 $U$ 的标准态射。于是序列
\begin{equation*}
0\to\cO_{\widehat{X},P}\otimes_{\cO_{X,P}}M'\to\cO_{\widehat{X},P}\otimes_{\cO_{X,P}}M\to\cO_{\widehat{X},P}\otimes_{\cO_{X,P}}M''\to 0
\end{equation*}
正合。因此 $\cO_{\widehat{X},P}$ 是平坦 $\cO_{X,P}$-代数。
\end{proof}

\begin{tnote}
此处「$i^*$ 正合」的实质是：形式完全化的茎 $\cO_{\widehat{X},P}$ 是 $\cO_{X,P}$ 的平坦代数。这正是形式函数定理的平坦版本，也是 2.7.8 能用 $i^*$ 比较 $\Ext$ 的原因；原书用 1.4.1 把有限生成模的局部正合性搬回概形上。[依据：原书 p.230–231]
\end{tnote}

\begin{proposition}\label{it:2.7.8}
设 $A$ 是诺特环，$I$ 是 $A$ 的理想，$X$ 是 $\Spec A$ 上的固有概形，$\cF$ 与 $\cG$ 是凝聚 $\cO_X$-模，$\widehat{X}$、$\widehat{\cF}$ 与 $\widehat{\cG}$ 分别是 $X$、$\cF$ 与 $\cG$ 关于 $I^{\sim}\cO_X$ 的形式完全化。
\begin{enumerate}[label=(\roman*),leftmargin=2.8em]
\item 对每个 $n$，$\Ext^n_{\cO_X}(\cF,\cG)$ 是凝聚 $\cO_X$-模，且
\begin{equation*}
(\Ext^n_{\cO_X}(\cF,\cG))^{\wedge}\cong\Ext^n_{\cO_{\widehat{X}}}(\widehat{\cF},\widehat{\cG}),
\end{equation*}
其中 $(\Ext^n_{\cO_X}(\cF,\cG))^{\wedge}$ 是 $\Ext^n_{\cO_X}(\cF,\cG)$ 关于 $I^{\sim}\cO_X$ 的形式完全化。
\item 对每个 $n$，$\Ext^n_{\cO_X}(\cF,\cG)$ 是有限生成 $A$-模，且
\begin{equation*}
(\Ext^n_{\cO_X}(\cF,\cG))^{\wedge}\cong\Ext^n_{\cO_{\widehat{X}}}(\widehat{\cF},\widehat{\cG}),
\end{equation*}
其中 $(\Ext^n_{\cO_X}(\cF,\cG))^{\wedge}$ 是 $\Ext^n_{\cO_X}(\cF,\cG)$ 的 $I$-进完全化。
\end{enumerate}
\end{proposition}

\begin{proof}
(i) 设
\begin{equation*}
0\to\cG\to\cI^{\bullet}
\end{equation*}
是 $\cO_X$-模 $\cG$ 的内射分解，并设
\begin{equation*}
0\to\widehat{\cG}\to\cJ^{\bullet}
\end{equation*}
是 $\cO_{\widehat{X}}$-模 $\widehat{\cG}$ 的内射分解。我们有
\begin{equation*}
\Ext^n_{\cO_X}(\cF,\cG)\cong H^n(\Hom_{\cO_X}(\cF,\cI^{\bullet})).
\end{equation*}
于是由 2.7.7 有
\begin{equation*}
i^*\Ext^n_{\cO_X}(\cF,\cG)\cong H^n(i^*\Hom_{\cO_X}(\cF,\cI^{\bullet})).
\end{equation*}
我们有标准同态
\begin{equation*}
i^*\Hom_{\cO_X}(\cF,\cI^{\bullet})\to\Hom_{\cO_{\widehat{X}}}(i^*\cF,i^*\cI^{\bullet}).
\end{equation*}
再次由 2.7.7，序列
\begin{equation*}
0\to i^*\cG\to i^*\cI^{\bullet}
\end{equation*}
正合。由 1.5.13 (ii) 有 $i^*\cG\cong\widehat{\cG}$。由 2.1.1 (ii)，这个同构可以延拓为复形的态射
\begin{equation*}
i^*\cI^{\bullet}\to\cJ^{\bullet}.
\end{equation*}
由 2.5.4 (i)，$\Ext^n_{\cO_X}(\cF,\cG)$ 是凝聚的。于是由 1.5.13 有
\begin{equation*}
(\Ext^n_{\cO_X}(\cF,\cG))^{\wedge}=i^*\Ext^n_{\cO_X}(\cF,\cG).
\end{equation*}
我们定义一个态射
\begin{equation*}
(\Ext^n_{\cO_X}(\cF,\cG))^{\wedge}\to\Ext^n_{\cO_{\widehat{X}}}(\widehat{\cF},\widehat{\cG})
\end{equation*}
为如下复合：
\begin{align*}
(\Ext^n_{\cO_X}(\cF,\cG))^{\wedge}&\cong i^*\Ext^n_{\cO_X}(\cF,\cG)\\
&\cong H^n(i^*\Hom_{\cO_X}(\cF,\cI^{\bullet}))\\
&\to H^n(\Hom_{\cO_{\widehat{X}}}(i^*\cF,i^*\cI^{\bullet}))\\
&\to H^n(\Hom_{\cO_{\widehat{X}}}(i^*\cF,\cJ^{\bullet}))\\
&\cong\Ext^n_{\cO_{\widehat{X}}}(\widehat{\cF},\widehat{\cG}).
\end{align*}
我们来证明它是同构。这是局部问题。于是我们可以假设存在 $\cO_X$-模的正合序列
\begin{equation*}
\cdots\to\cE_1\to\cE_0\to\cF\to 0,
\end{equation*}
其中每个 $\cE_n$ 是有限秩自由模。由 2.7.7，序列
\begin{equation*}
\cdots\to i^*\cE_1\to i^*\cE_0\to i^*\cF\to 0
\end{equation*}
正合，且每个 $i^*\cE_n$ 是有限秩自由的。由 2.2.16，我们有
\begin{align*}
\Ext^n_{\cO_X}(\cF,\cG)&\cong H^n(\Hom_{\cO_X}(\cE_{\bullet},\cG)),\\
\Ext^n_{\cO_{\widehat{X}}}(\widehat{\cF},\widehat{\cG})&\cong H^n(\Hom_{\cO_{\widehat{X}}}(i^*\cE_{\bullet},i^*\cG)).
\end{align*}
由 2.7.7 有
\begin{equation*}
(\Ext^n_{\cO_X}(\cF,\cG))^{\wedge}\cong i^*\Ext^n_{\cO_X}(\cF,\cG)\cong H^n(i^*\Hom_{\cO_X}(\cE_{\bullet},\cG)).
\end{equation*}
由于 $\cE_{\bullet}$ 是自由 $\cO_X$-模的复形，我们有
\begin{equation*}
i^*\Hom_{\cO_X}(\cE_{\bullet},\cG)\cong\Hom_{\cO_{\widehat{X}}}(i^*\cE_{\bullet},i^*\cG).
\end{equation*}
把所有这些同构放在一起，我们得到
\begin{equation*}
(\Ext^n_{\cO_X}(\cF,\cG))^{\wedge}\cong\Ext^n_{\cO_{\widehat{X}}}(\widehat{\cF},\widehat{\cG}).
\end{equation*}
(ii) 我们有双正则谱序列
\begin{equation*}
H^p(X,\Ext^q_{\cO_X}(\cF,\cG))\Rightarrow\Ext^{p+q}(\cF,\cG).
\end{equation*}
由 2.5.4 (ii)，这个谱序列只涉及有限生成 $A$-模。取 $I$-进完全化并应用 1.5.4，我们得到双正则谱序列
\begin{equation*}\tag{1}
(H^p(X,\Ext^q_{\cO_X}(\cF,\cG)))^{\wedge}\Rightarrow(\Ext^{p+q}_{\cO_X}(\cF,\cG))^{\wedge}.
\end{equation*}
另一方面，我们有双正则谱序列
\begin{equation*}\tag{2}
H^p(\widehat{X},\Ext^q_{\cO_{\widehat{X}}}(\widehat{\cF},\widehat{\cG}))\Rightarrow\Ext^{p+q}_{\cO_{\widehat{X}}}(\widehat{\cF},\widehat{\cG}).
\end{equation*}
我们留给读者去构造从谱序列 (1) 到谱序列 (2) 的态射。由 2.7.5 我们有
\begin{equation*}
(H^p(X,\Ext^q_{\cO_X}(\cF,\cG)))^{\wedge}\cong H^p(\widehat{X},(\Ext^q_{\cO_X}(\cF,\cG))^{\wedge})\cong H^p(\widehat{X},\Ext^q_{\cO_{\widehat{X}}}(\widehat{\cF},\widehat{\cG})).
\end{equation*}
于是由 2.2.1 有
\begin{equation*}
(\Ext^n_{\cO_X}(\cF,\cG))^{\wedge}\cong\Ext^n_{\cO_{\widehat{X}}}(\widehat{\cF},\widehat{\cG}).
\end{equation*}
\end{proof}

\begin{corollary}\label{it:2.7.9}
保持 2.7.8 的假设与记号。若 $A$ 是 $I$-进完备的，则我们有
\begin{equation*}
\Hom_{\cO_X}(\cF,\cG)\cong\Hom_{\cO_{\widehat{X}}}(\widehat{\cF},\widehat{\cG}),
\end{equation*}
其中该同构把每个态射 $u\colon\cF\to\cG$ 映为它在形式完全化上诱导的态射 $\hat{u}\colon\widehat{\cF}\to\widehat{\cG}$。此外，$u$ 是单态射（resp.\ 满态射、同构）当且仅当 $\hat{u}$ 是单态射（resp.\ 满态射、同构）。
\end{corollary}

\begin{proof}
在 2.7.8 (ii) 中取 $n=0$，我们得到
\begin{equation*}
(\Hom_{\cO_X}(\cF,\cG))^{\wedge}\cong\Hom_{\cO_{\widehat{X}}}(\widehat{\cF},\widehat{\cG}).
\end{equation*}
由 2.5.4 (ii)，$\Hom_{\cO_X}(\cF,\cG)$ 是有限生成 $A$-模。由于 $A$ 是 $I$-进完备的，由 1.5.5 我们有
\begin{equation*}
\Hom_{\cO_X}(\cF,\cG)\cong(\Hom_{\cO_X}(\cF,\cG))^{\wedge}.
\end{equation*}
因此
\begin{equation*}
\Hom_{\cO_X}(\cF,\cG)\cong\Hom_{\cO_{\widehat{X}}}(\widehat{\cF},\widehat{\cG}).
\end{equation*}
由 1.5.16 (ii)，态射 $\hat{u}\colon\widehat{\cF}\to\widehat{\cG}$ 是单态射（resp.\ 满态射、同构）当且仅当 $u\colon\cF\to\cG$ 在 $X$ 中 $\widehat{X}$ 的某个邻域内是单态射（resp.\ 满态射、同构）。我们断言 $X$ 是 $\widehat{X}$ 在 $X$ 中的唯一邻域。这就证明了我们的断言。设 $U$ 是 $\widehat{X}$ 在 $X$ 中的开邻域。由于 $f$ 是固有的，$f(X-U)$ 是 $\Spec A$ 的闭子集。于是存在 $A$ 的理想 $\mathfrak{a}$，使得 $f(X-U)=V(\mathfrak{a})$。另一方面，我们有 $\widehat{X}=f^{-1}(V(I))$。于是 $f(X-\widehat{X})\cap V(I)=\varnothing$。特别地，我们有 $f(X-U)\cap V(I)=\varnothing$，即 $V(\mathfrak{a})\cap V(I)=\varnothing$。由于 $A$ 是 $I$-进完备的，由 1.5.7 有 $I\subset\operatorname{rad}(A)$，从而 $A$ 的所有极大理想都含于 $V(I)$。于是没有 $A$ 的极大理想含于 $V(\mathfrak{a})$。这只可能在 $\mathfrak{a}=A$ 时发生，即 $V(\mathfrak{a})=\varnothing$。所以 $f(X-U)=\varnothing$，从而 $X=U$。
\end{proof}

\begin{theorem}[Grothendieck 存在定理]\label{it:2.7.10}
设 $A$ 是诺特环，$I$ 是 $A$ 的理想，$X$ 是 $\Spec A$ 上的固有概形，$\widehat{X}$ 是 $X$ 关于理想层 $I^{\sim}\cO_X$ 的形式完全化。若 $A$ 是 $I$-进完备的，则对每个凝聚 $\cO_{\widehat{X}}$-模 $\fF$，存在凝聚 $\cO_X$-模 $\cF$，它在唯一同构意义下唯一，使得 $\fF\cong\widehat{\cF}$，其中 $\widehat{\cF}$ 是 $\cF$ 关于理想层 $I^{\sim}\cO_X$ 的形式完全化。
\end{theorem}

考虑到凝聚 $\cO_{\widehat{X}}$-模的定义、1.5.13 (i) 与 1.5.18，该定理可以改述为如下不涉及形式完全化的形式：

\begin{theorem}[Grothendieck 存在定理]\label{it:2.7.11}
设 $A$ 是诺特环，$I$ 是 $A$ 的理想，$X$ 是 $\Spec A$ 上的固有概形。对每个 $k$，设 $X_k$ 是 $X$ 的由理想层 $(I^k)^{\sim}\cO_X$ 定义的闭子概形。用如下交换图固定记号（$l\ge k$）：

\begin{center}
\begin{tikzcd}
(X_k,\cO_{X_k}) \arrow[r,"i_k"] \arrow[d,"i_{kl}"'] & (X,\cO_X) \\
(X_l,\cO_{X_l}) \arrow[r,"\hat{i}_l"] \arrow[ur,"i_l"'] & (\widehat{X},\cO_{\widehat{X}}) \arrow[u,"i"']
\end{tikzcd}
\end{center}

其中所有箭头都是标准态射。设 $(F_k,\phi_{lk})_{k\in\bN}$ 是由对每个 $k$ 的一个 $\cO_{X_k}$-模 $F_k$ 与对每对 $l\ge k$ 的一个态射 $\phi_{lk}\colon F_l\to i_{kl*}F_k$ 组成的族，使得下列条件成立：
\begin{enumerate}[label=(\alph*),leftmargin=2.8em]
\item 每个 $F_k$ 是凝聚 $\cO_{X_k}$-模。
\item 对每个 $k$ 有 $\phi_{kk}=\id_{F_k}$，且对每个三元组 $m\ge l\ge k$ 有 $i_{lm*}(\phi_{lk})\circ\phi_{ml}=\phi_{mk}$。
\item 对每对 $l\ge k$，$\phi_{lk}$ 诱导同构 $i_{kl}^*F_l\cong F_k$，即 $F_l\otimes_{\cO_{X_l}}\cO_{X_k}\cong F_k$。
\end{enumerate}
若 $A$ 是 $I$-进完备的，则存在凝聚 $\cO_X$-模 $\cF$，它在唯一同构意义下唯一，使得族 $(F_k,\phi_{lk})_{k\in\bN}$ 同构于族 $(i_k^*\cF,\psi_{lk})_{k\in\bN}$，其中对每对 $l\ge k$，$\psi_{lk}$ 是标准态射
\begin{equation*}
i_l^*\cF\to i_{kl*}i_{kl}^*i_l^*\cF\cong i_{kl*}i_k^*\cF.
\end{equation*}
\end{theorem}

\begin{tnote}
2.7.11 与 2.7.10 是同一命题的两种形式：前者只用闭子概形 $X_k$ 与相容族陈述，避免提到形式完全化。注意原书随后实际证明的是 2.7.10（见 p.242 「Finally, we are ready to prove 2.7.10」）；2.7.11 的存在性部分由唯一性论证与 1.5.13、1.5.18 提供的字典得到。[依据：原书 p.234、p.242]
\end{tnote}

该定理的唯一性部分直接由 2.7.9 得到。在证明定理的存在性部分之前，我们需要一系列引理。回忆：凝聚 $\cO_{\widehat{X}}$-模 $\fF$ 称为\emph{可代数化}的，若存在凝聚 $\cO_X$-模 $\cF$ 使得 $\fF\cong\widehat{\cF}$。以下我们在 2.7.10 的假设下工作。

\begin{lemma}\label{it:2.7.12}
设 $\fF$ 与 $\fG$ 是两个凝聚 $\cO_{\widehat{X}}$-模，$u\colon\fF\to\fG$ 是态射。若 $\fF$ 与 $\fG$ 可代数化，则 $\ker u$、$\coker u$ 与 $\Image u$ 都可代数化。
\end{lemma}

\begin{proof}
设 $\cF$ 与 $\cG$ 是凝聚 $\cO_X$-模，使得 $\fF\cong\widehat{\cF}$，$\fG\cong\widehat{\cG}$。由 2.7.9，存在态射 $u'\colon\cF\to\cG$，使得 $u=\widehat{u'}$。由 1.5.13 (i)，我们有 $\ker u\cong\widehat{\ker u'}$。故 $\ker u$ 可代数化。类似地，$\coker u$ 与 $\Image u$ 也可代数化。
\end{proof}

\begin{lemma}\label{it:2.7.13}
设
\begin{equation*}
0\to\fF\to\fG\to\fH\to 0
\end{equation*}
是 $\cO_{\widehat{X}}$-模的正合序列。假设 $\fF$ 与 $\fH$ 是可代数化的凝聚 $\cO_{\widehat{X}}$-模。则 $\fG$ 也是可代数化的凝聚 $\cO_{\widehat{X}}$-模。
\end{lemma}

\begin{proof}
设 $\cF$ 与 $\cH$ 是凝聚 $\cO_X$-模，使得 $\fF\cong\widehat{\cF}$，$\fH\cong\widehat{\cH}$。由 2.7.8 (ii) 我们有
\begin{equation*}
(\Ext^1_{\cO_X}(\cH,\cF))^{\wedge}\cong\Ext^1_{\cO_{\widehat{X}}}(\fH,\fF).
\end{equation*}
由 2.5.4 (ii)，$\Ext^1_{\cO_X}(\cH,\cF)$ 是有限生成 $A$-模。由于 $A$ 是 $I$-进完备的，由 1.5.5 有
\begin{equation*}
(\Ext^1_{\cO_X}(\cH,\cF))^{\wedge}\cong\Ext^1_{\cO_X}(\cH,\cF).
\end{equation*}
所以
\begin{equation*}
\Ext^1_{\cO_X}(\cH,\cF)\cong\Ext^1_{\cO_{\widehat{X}}}(\fH,\fF).
\end{equation*}
由 2.1.15，正合序列
\begin{equation*}
0\to\fF\to\fG\to\fH\to 0
\end{equation*}
对应于 $\Ext^1_{\cO_{\widehat{X}}}(\fH,\fF)$ 中的一个元素，从而对应于 $\Ext^1_{\cO_X}(\cH,\cF)$ 中的一个元素。设
\begin{equation*}
0\to\cF\to\cG\to\cH\to 0
\end{equation*}
是对应于该元素的短正合序列。由 2.4.7，$\cG$ 是凝聚 $\cO_X$-模。此外我们有 $\fG\cong\widehat{\cG}$。所以 $\fG$ 是可代数化的凝聚 $\cO_{\widehat{X}}$-模。
\end{proof}
